\section{问题重述}
\subsection{问题背景}
题目给出的通用神经网络处理器在多个核心上执行同一张计算图。矩阵、向量及搬运操作之间有先后依赖；每核片上存储有限，所有核心共享外部存储带宽。切图既释放并行性，也改变跨核搬运、等待和张量驻留时间，因此需要同时确定操作归属与执行顺序，并以总体任务执行时间（Makespan）评价方案。

本届A题提供计算图、硬件配置和评估程序。参赛者提交全部非COPY操作的子图归属及各核的子图顺序；评估程序据此生成搬运与缓存换出换入操作，计算Makespan、额外数据搬运量和缓存命中率。L1和UB由各核独占，DDR由多核共享，问题三另设共享只读L2 Cache。

\subsection{问题描述}
在相同计算图和固定硬件配置下，题目依次改变Task组织、片上复用和共享读服务条件。各问均须给出可执行的多核计划及相应评价指标。

\textbf{问题一：}场景A将每个子图构造成独立Task，跨子图张量通过DDR中转，同核相邻Task也不保留片上副本。需要确定全部非COPY操作的子图归属和每核的子图执行次序，在满足计算依赖与Task等待规则的前提下缩短总体工期、控制额外搬运，并给出1至5核的加速比及逐图结果。

\textbf{问题二：}场景B把同核子图合并为核级Task，使同核依赖可以复用L1和UB中的数据；跨核依赖仍需写出、读入和固定同步延迟。需要在容量约束下重新设计切图、分核和核内顺序，给出不同核数的工期、加速比与额外搬运，并分析同核复用带来的收益和换出恢复的代价。

\textbf{问题三：}在问题二的核级Task结构上增加共享只读L2 Cache，本文将此配置记为C。读请求在发射时查询，完成后尝试插入，空间不足时按FIFO淘汰。需比较同图同核数下有无L2的Makespan与额外数据搬运量，给出命中率，并设计适应共享读资源的方案。
